% ============================================================
%  EE3165 -- IC Design  |  Laboratory 1
%  Experiment 3 & 4: FSM Implementation
%  Compiler : pdfLaTeX (Overleaf default)
%  Encoding : utf8 + T5 + babel(vietnamese)
% ============================================================
\documentclass[12pt, a4paper]{article}

% ---- Encoding & Language -----------------------------------
\usepackage[utf8]{inputenc}
\usepackage[T5]{fontenc}
\usepackage[vietnamese]{babel}

% ---- Page layout -------------------------------------------
\usepackage[
    a4paper,
    top=2.5cm, bottom=2.5cm,
    left=3.0cm, right=2.2cm,
    headheight=15pt, headsep=0.8cm, footskip=1cm
]{geometry}

% ---- Header / Footer ---------------------------------------
\usepackage{fancyhdr}
\pagestyle{fancy}
\fancyhf{}
\fancyhead[L]{\small\itshape EE3165 -- IC Design}
\fancyhead[R]{\small\itshape Laboratory 1: FSM Implementation}
\fancyfoot[C]{\small\thepage}
\renewcommand{\headrulewidth}{0.4pt}
\renewcommand{\footrulewidth}{0.4pt}

% ---- Math --------------------------------------------------
\usepackage{amsmath, amssymb}

% ---- Code listing ------------------------------------------
\usepackage{listings}
\usepackage[dvipsnames]{xcolor}

\definecolor{codebg}{RGB}{248,248,248}
\definecolor{codeframe}{RGB}{200,200,200}
\definecolor{kwcolor}{RGB}{0,0,180}
\definecolor{cmtcolor}{RGB}{60,128,49}
\definecolor{numcolor}{RGB}{130,130,130}

\lstdefinelanguage{SystemVerilog}{
    morekeywords={module,endmodule,input,output,logic,wire,reg,
                  always_ff,always_comb,always,assign,begin,end,
                  if,else,case,endcase,posedge,negedge,typedef,
                  enum,parameter,localparam,for,initial,timescale,
                  default,integer,task,endtask,posedge,negedge},
    sensitive=true,
    morecomment=[l]{//},
    morecomment=[s]{/*}{*/},
    morestring=[b]"
}

\lstset{
    language=SystemVerilog,
    basicstyle=\ttfamily\footnotesize,
    keywordstyle=\color{kwcolor}\bfseries,
    commentstyle=\color{cmtcolor}\itshape,
    backgroundcolor=\color{codebg},
    frame=single,
    rulecolor=\color{codeframe},
    numbers=left,
    numberstyle=\tiny\color{numcolor},
    numbersep=8pt,
    breaklines=true,
    showstringspaces=false,
    tabsize=4,
    captionpos=b,
    xleftmargin=2em,
    framexleftmargin=1.8em,
    aboveskip=8pt,
    belowskip=4pt
}

% Bash/shell style for filelist
\lstdefinestyle{bash}{
    language=bash,
    basicstyle=\ttfamily\footnotesize,
    backgroundcolor=\color{black!90},
    keywordstyle=\color{green!70},
    commentstyle=\color{gray},
    frame=single,
    rulecolor=\color{black},
    numbers=none,
    breaklines=true,
    showstringspaces=false,
    aboveskip=8pt,
    belowskip=4pt,
    xleftmargin=1em,
    framexleftmargin=0.8em
}

% Plain text style for filelist.f
\lstdefinestyle{plainfile}{
    basicstyle=\ttfamily\footnotesize\color{black},
    backgroundcolor=\color{gray!8},
    frame=single,
    rulecolor=\color{codeframe},
    numbers=none,
    breaklines=true,
    aboveskip=6pt,
    belowskip=4pt,
    xleftmargin=1em,
    framexleftmargin=0.8em
}

% ---- Tables ------------------------------------------------
\usepackage{booktabs}
\usepackage{array}
\usepackage{multirow}
\renewcommand{\arraystretch}{1.35}

% ---- Graphics ----------------------------------------------
\usepackage{graphicx}
\usepackage{float}
\usepackage{caption}
\captionsetup{font=small, labelfont=bf, labelsep=period,
              justification=centering, skip=6pt}

% ---- Boxes -------------------------------------------------
\usepackage{mdframed}
\newmdenv[linecolor=blue!50!black,linewidth=1pt,
    leftline=true,rightline=false,topline=false,bottomline=false,
    backgroundcolor=blue!4,
    innerleftmargin=10pt,innerrightmargin=8pt,
    innertopmargin=5pt,innerbottommargin=5pt,
    skipabove=8pt,skipbelow=8pt]{infobox}
\newmdenv[linecolor=OliveGreen,linewidth=1pt,
    leftline=true,rightline=false,topline=false,bottomline=false,
    backgroundcolor=green!4,
    innerleftmargin=10pt,innerrightmargin=8pt,
    innertopmargin=5pt,innerbottommargin=5pt,
    skipabove=8pt,skipbelow=8pt]{resultbox}
\newmdenv[linecolor=orange!70!black,linewidth=1pt,
    leftline=true,rightline=false,topline=false,bottomline=false,
    backgroundcolor=orange!5,
    innerleftmargin=10pt,innerrightmargin=8pt,
    innertopmargin=5pt,innerbottommargin=5pt,
    skipabove=8pt,skipbelow=8pt]{warnbox}

% ---- Lists -------------------------------------------------
\usepackage{enumitem}
\setlist[itemize]{leftmargin=1.5cm,itemsep=2pt,topsep=3pt}
\setlist[enumerate]{leftmargin=1.5cm,itemsep=2pt,topsep=3pt}

% ---- Section formatting ------------------------------------
\usepackage{titlesec}
\titleformat{\section}{\large\bfseries}{\thesection.}{0.5em}{}
    [\vspace{-4pt}\rule{\textwidth}{0.4pt}]
\titleformat{\subsection}{\normalsize\bfseries}{\thesubsection.}{0.5em}{}
\titleformat{\subsubsection}{\normalsize\itshape\bfseries}{\thesubsubsection.}{0.5em}{}
\titlespacing*{\section}{0pt}{16pt}{8pt}
\titlespacing*{\subsection}{0pt}{10pt}{4pt}
\titlespacing*{\subsubsection}{0pt}{8pt}{3pt}

% ---- Misc --------------------------------------------------
\usepackage[nottoc]{tocbibind}
\usepackage[colorlinks=true,linkcolor=blue!60!black,
            urlcolor=blue!70!black,bookmarksnumbered]{hyperref}
\usepackage{parskip}
\setlength{\parskip}{5pt}
\setlength{\parindent}{1cm}
\usepackage{indentfirst}
\usepackage{setspace}
\onehalfspacing
\counterwithout{figure}{section}
\counterwithout{table}{section}

\usepackage{xcolor}
\usepackage{listings}

\definecolor{codegray}{rgb}{0.2,0.2,0.2}
\definecolor{codewhite}{rgb}{1,1,1}

\lstdefinestyle{bash_custom}{
    backgroundcolor=\color{codegray},   % màu nền xám đậm
    basicstyle=\ttfamily\small\color{codewhite}, % màu chữ trắng
    breaklines=true,
    frame=single,                       % tạo khung nếu cần
    rulecolor=\color{white},            % màu của khung
    showstringspaces=false
}

% ============================================================
\begin{document}
% ============================================================

% ---- Cover -------------------------------------------------
\begin{center}
    {\large\bfseries TRƯỜNG ĐẠI HỌC BÁCH KHOA TP.~HỒ CHÍ MINH}\\[4pt]
    {\normalsize Khoa Điện tử -- Viễn thông}\\[0.7cm]
    \rule{\textwidth}{1.5pt}\\[6pt]
    {\LARGE\bfseries LABORATORY 1}\\[4pt]
    {\large Hardware Description and Digital Design Verification}\\[3pt]
    {\normalsize\itshape IC Design (EE3165)}\\[6pt]
    \rule{\textwidth}{0.8pt}\\[0.5cm]
    {\large\bfseries Experiment 3 \& 4: FSM Implementation}\\[1cm]
    \begin{tabular}{ll}
        \textbf{Họ và tên:} & Thái Đức Thiên \\[5pt]
        \textbf{MSSV:}      & 2313222 \\[5pt]
        \textbf{Lớp:}       & L03 \\[5pt]
        \textbf{Ngày nộp:}  & \the\day/\the\month/\the\year \\
    \end{tabular}
\end{center}

\newpage
\tableofcontents
\newpage

% ============================================================
\section{Experiment 3: FSM -- Logic-Expression Style}
\label{sec:exp3}
% ============================================================

\subsection{Mô tả bài toán}
\label{subsec:exp3-desc}

Thiết kế một FSM nhận diện hai chuỗi ký hiệu: \textbf{bốn bit 1 liên tiếp} hoặc \textbf{bốn bit 0 liên tiếp}. Hệ thống có một đầu vào $w$ và một đầu ra $z$. Khi $w$ duy trì giá trị không đổi trong bốn chu kỳ clock liên tiếp, $z = 1$; ngược lại $z = 0$. Overlapping sequence được cho phép.

\subsection{Phân tích FSM}
\label{subsec:exp3-fsm}

FSM gồm \textbf{9 trạng thái} (A đến I). Trạng thái ban đầu (sau reset) là \textbf{A}. Bảng chuyển trạng thái được rút ra từ Figure~4:

\begin{table}[H]
\centering
\caption{Bảng chuyển trạng thái FSM.}
\label{tab:state-transition}
\begin{tabular}{ccccc}
\toprule
\textbf{State hiện tại} & \textbf{$w=0$ (next)} & \textbf{$w=1$ (next)} & \textbf{$z$} & \textbf{Ý nghĩa} \\
\midrule
A & B & F & 0 & Reset / khởi đầu \\
B & C & F & 0 & Đếm 1 bit 0 \\
C & D & F & 0 & Đếm 2 bit 0 \\
D & E & F & 0 & Đếm 3 bit 0 \\
E & E & F & 1 & $\geq$4 bit 0 (overlapping) \\
F & B & G & 0 & Đếm 1 bit 1 \\
G & B & H & 0 & Đếm 2 bit 1 \\
H & B & I & 0 & Đếm 3 bit 1 \\
I & B & I & 1 & $\geq$4 bit 1 (overlapping) \\
\bottomrule
\end{tabular}
\end{table}

Đầu ra $z = 1$ chỉ tại trạng thái \textbf{E} (phát hiện 4 bit 0) và \textbf{I} (phát hiện 4 bit 1).

\subsection{Cấu trúc thư mục dự án}
\label{subsec:exp3-tree}

\begin{warnbox}
Cấu trúc thư mục dưới đây tuân theo quy ước của môn học: \texttt{00\_src} chứa các file RTL design, \texttt{01\_tb} chứa testbench, và \texttt{filelist.f} liệt kê toàn bộ file cần biên dịch.
\end{warnbox}

\begin{lstlisting}[style=bash_custom, caption={Cấu trúc thư mục Experiment 3 (lệnh \texttt{tree} trên Linux).}]
Lab1_xcelium/
|-- 00_src/
|   |-- fsm_binary.sv
|   `-- fsm_onehot.sv
|-- 01_tb/
|   |-- tb_fsm_binary.sv
|   `-- tb_fsm_onehot.sv
`-- filelist.f
\end{lstlisting}

\subsection{File \texttt{filelist.f}}
\label{subsec:exp3-filelist}

File \texttt{filelist.f} liệt kê đường dẫn tương đối đến tất cả các file SystemVerilog cần biên dịch. Khi chạy Xcelium, lệnh \texttt{xrun -f filelist.f} sẽ đọc file này và biên dịch toàn bộ theo thứ tự.

\begin{lstlisting}[style=plainfile, caption={Nội dung file \texttt{filelist.f} -- Experiment 3.}]
./00_src/fsm_binary.sv
./00_src/fsm_onehot.sv
./01_tb/tb_fsm_binary.sv
./01_tb/tb_fsm_onehot.sv
\end{lstlisting}

\begin{infobox}
\textbf{Lệnh chạy mô phỏng trên Xcelium (Linux terminal):}\\[4pt]
\texttt{xrun -f filelist.f -top tb\_fsm\_binary -access +rwc}\\
\texttt{xrun -f filelist.f -top tb\_fsm\_onehot -access +rwc}
\end{infobox}

\subsection{File \texttt{00\_src} -- RTL Design}
\label{subsec:exp3-src}

\subsubsection{Binary Encoding: \texttt{fsm\_binary.sv}}

Phong cách \textbf{logic-expression}: next-state được mô tả bằng chuỗi toán tử điều kiện (\texttt{?:}) nối tiếp nhau trực tiếp trong \texttt{assign} statement, không dùng \texttt{always\_comb}. Đây là cách tiếp cận gần với mức cổng logic nhất.

\begin{lstlisting}[caption={\texttt{00\_src/fsm\_binary.sv} -- Binary Encoding, Logic-Expression Style.},
                   label={lst:exp3-binary}]
`timescale 1ns/1ns
module fsm_binary (
    input  logic clk,
    input  logic rst,
    input  logic w,
    output logic z
);

    logic [3:0] state, next_state;

    localparam logic [3:0]
        A = 4'b0000,
        B = 4'b0001,
        C = 4'b0010,
        D = 4'b0011,
        E = 4'b0100,
        F = 4'b0101,
        G = 4'b0110,
        H = 4'b0111,
        I = 4'b1000;

    // State register (synchronous reset)
    always_ff @(posedge clk) begin
        if (rst)
            state <= A;
        else
            state <= next_state;
    end

    // Next-state logic (logic-expression / conditional chain)
    assign next_state =
        (state == A) ? (w ? F : B) :
        (state == B) ? (w ? F : C) :
        (state == C) ? (w ? F : D) :
        (state == D) ? (w ? F : E) :
        (state == E) ? (w ? F : E) :
        (state == F) ? (w ? G : B) :
        (state == G) ? (w ? H : B) :
        (state == H) ? (w ? I : B) :
        (state == I) ? (w ? I : B) :
        A;

    // Output logic: z=1 in states E (four 0s) and I (four 1s)
    assign z = (state == E) || (state == I);

endmodule
\end{lstlisting}

\subsubsection{One-Hot Encoding: \texttt{fsm\_onehot.sv}}

Phong cách \textbf{logic-expression}: mỗi bit của \texttt{next\_state} được mô tả bằng một \texttt{assign} statement độc lập, tương ứng trực tiếp với các phương trình Boolean rút ra từ bảng chuyển trạng thái.

Phương trình next-state theo one-hot (với \texttt{state[k]} là flip-flop của trạng thái thứ $k$):

\begin{align}
    \text{next}[0] &= 0 \quad\text{(không state nào chuyển về A, chỉ reset)} \notag\\
    \text{next}[1] &= \bar{w}\cdot(\text{state}[0] \,|\, \text{state}[5] \,|\, \text{state}[6] \,|\, \text{state}[7] \,|\, \text{state}[8]) \notag\\
    \text{next}[2] &= \bar{w}\cdot\text{state}[1] \notag\\
    \text{next}[3] &= \bar{w}\cdot\text{state}[2] \notag\\
    \text{next}[4] &= \bar{w}\cdot(\text{state}[3] \,|\, \text{state}[4]) \quad\text{(overlapping tại E)} \notag\\
    \text{next}[5] &= w\cdot(\text{state}[0] \,|\, \text{state}[1] \,|\, \text{state}[2] \,|\, \text{state}[3] \,|\, \text{state}[4]) \notag\\
    \text{next}[6] &= w\cdot\text{state}[5] \notag\\
    \text{next}[7] &= w\cdot\text{state}[6] \notag\\
    \text{next}[8] &= w\cdot(\text{state}[7] \,|\, \text{state}[8]) \quad\text{(overlapping tại I)} \notag
\end{align}

\begin{lstlisting}[caption={\texttt{00\_src/fsm\_onehot.sv} -- One-Hot Encoding, Logic-Expression Style.},
                   label={lst:exp3-onehot}]
`timescale 1ns/1ns
module fsm_onehot (
    input  logic clk,
    input  logic rst,
    input  logic w,
    output logic z
);

    // state[0]=A, [1]=B, [2]=C, [3]=D, [4]=E
    // state[5]=F, [6]=G, [7]=H, [8]=I
    logic [8:0] state, next_state;

    // State register (synchronous reset -> state A = 9'b000000001)
    always_ff @(posedge clk) begin
        if (rst)
            state <= 9'b000000001;
        else
            state <= next_state;
    end

    // Next-state logic (direct Boolean equations from state diagram)
    assign next_state[0] = 1'b0;
    assign next_state[1] = ~w & (state[0] | state[5] | state[6] | state[7] | state[8]);
    assign next_state[2] = ~w & state[1];
    assign next_state[3] = ~w & state[2];
    assign next_state[4] = ~w & (state[3] | state[4]);   // overlapping at E
    assign next_state[5] =  w & (state[0] | state[1] | state[2] | state[3] | state[4]);
    assign next_state[6] =  w & state[5];
    assign next_state[7] =  w & state[6];
    assign next_state[8] =  w & (state[7] | state[8]);   // overlapping at I

    // Output: z=1 in state E (bit 4) or state I (bit 8)
    assign z = state[4] | state[8];

endmodule
\end{lstlisting}

\subsection{File \texttt{01\_tb} -- Testbench}
\label{subsec:exp3-tb}

\subsubsection{Testbench One-Hot: \texttt{tb\_fsm\_onehot.sv}}

\begin{lstlisting}[caption={\texttt{01\_tb/tb\_fsm\_onehot.sv} -- Testbench cho fsm\_onehot.},
                   label={lst:exp3-tb-oh}]
`timescale 1ns/1ns
module tb_fsm_onehot;
    logic clk;
    logic rst;
    logic w;
    logic z;

    fsm_onehot dut (.clk(clk), .rst(rst), .w(w), .z(z));

    // Clock generation: period = 10ns
    initial clk = 0;
    always #5 clk = ~clk;

    // Task: apply one bit to input w and advance one clock cycle
    task apply_bit(input logic x);
    begin
        w = x;
        @(posedge clk); #1;
    end
    endtask

    initial begin
        $dumpfile("tb_fsm_onehot.vcd");
        $dumpvars(0, tb_fsm_onehot);

        // Reset
        rst = 1;
        w   = 0;
        @(posedge clk); #1;
        rst = 0;

        // --- Test Case 1: four consecutive 0s -> z=1 on 4th cycle ---
        apply_bit(0);   // A->B
        apply_bit(0);   // B->C
        apply_bit(0);   // C->D
        apply_bit(0);   // D->E  (z=1)

        // --- Overlapping 0 (stay in E) ---
        apply_bit(0);   // E->E  (z=1)

        // --- Test Case 2: four consecutive 1s -> z=1 on 4th cycle ---
        apply_bit(1);   // E->F
        apply_bit(1);   // F->G
        apply_bit(1);   // G->H
        apply_bit(1);   // H->I  (z=1)

        // --- Overlapping 1 (stay in I) ---
        apply_bit(1);   // I->I  (z=1)

        // --- Test Case 3: mixed sequence (broken chain) ---
        apply_bit(0);   // I->B
        apply_bit(1);   // B->F
        apply_bit(0);   // F->B
        apply_bit(0);   // B->C
        apply_bit(0);   // C->D
        apply_bit(0);   // D->E  (z=1)

        #20;
        $finish;
    end
endmodule
\end{lstlisting}

\subsubsection{Testbench Binary: \texttt{tb\_fsm\_binary.sv}}

\begin{lstlisting}[caption={\texttt{01\_tb/tb\_fsm\_binary.sv} -- Testbench cho fsm\_binary.},
                   label={lst:exp3-tb-bin}]
`timescale 1ns/1ns
module tb_fsm_binary;
    logic clk;
    logic rst;
    logic w;
    logic z;

    fsm_binary dut (
        .clk(clk),
        .rst(rst),
        .w(w),
        .z(z)
    );

    // Clock generation
    initial clk = 1'b0;
    always #5 clk = ~clk;

    // Task: apply one bit and advance one clock cycle
    task apply_bit(input logic x);
    begin
        w = x;
        @(posedge clk); #1;
    end
    endtask

    initial begin
        $dumpfile("tb_fsm_binary.vcd");
        $dumpvars(0, tb_fsm_binary);

        // Reset (synchronous)
        rst = 1'b1;
        w   = 1'b0;
        @(posedge clk); #1;
        rst = 1'b0;

        // --- Test Case 1: four 0s -> z=1 ---
        apply_bit(0);   // A->B
        apply_bit(0);   // B->C
        apply_bit(0);   // C->D
        apply_bit(0);   // D->E  (z=1)

        // --- Overlapping 0 ---
        apply_bit(0);   // E->E  (z=1)

        // --- Test Case 2: four 1s -> z=1 ---
        apply_bit(1);   // E->F
        apply_bit(1);   // F->G
        apply_bit(1);   // G->H
        apply_bit(1);   // H->I  (z=1)

        // --- Overlapping 1 ---
        apply_bit(1);   // I->I  (z=1)

        // --- Test Case 3: mixed / broken chain ---
        apply_bit(0);   // I->B
        apply_bit(0);   // B->C
        apply_bit(1);   // C->F  (chain broken)
        apply_bit(0);   // F->B
        apply_bit(0);   // B->C
        apply_bit(0);   // C->D
        apply_bit(0);   // D->E  (z=1)

        // --- Test Case 4: reset mid-sequence ---
        rst = 1'b1;
        @(posedge clk); #1;
        rst = 1'b0;

        // Count 4 ones again after reset
        apply_bit(1);   // A->F
        apply_bit(1);   // F->G
        apply_bit(1);   // G->H
        apply_bit(1);   // H->I  (z=1)

        #20;
        $finish;
    end
endmodule
\end{lstlisting}

\subsection{Giải thích Waveform -- Experiment 3}
\label{subsec:exp3-wave}

% ----------------------------------------------------------
% [VỊ TRÍ HÌNH 1] Waveform Experiment 3 -- One-Hot
% Khi có ảnh: xóa \fbox{...}, bỏ comment \includegraphics
% \includegraphics[width=\textwidth]{wave_exp3_onehot.png}
% ----------------------------------------------------------
\begin{figure}[H]
    \centering
    \fbox{\parbox{0.95\textwidth}{\centering\vspace{2.2cm}
    \textit{[Vị trí Hình 1: Waveform Experiment 3 -- One-Hot Encoding (SimVision)]}
    \vspace{2.2cm}}}
    \caption{Waveform mô phỏng FSM Experiment 3 -- One-Hot Encoding.}
    \label{fig:wave-exp3-oh}
\end{figure}

% ----------------------------------------------------------
% [VỊ TRÍ HÌNH 2] Waveform Experiment 3 -- Binary
% \includegraphics[width=\textwidth]{wave_exp3_binary.png}
% ----------------------------------------------------------
\begin{figure}[H]
    \centering
    \fbox{\parbox{0.95\textwidth}{\centering\vspace{2.2cm}
    \textit{[Vị trí Hình 2: Waveform Experiment 3 -- Binary Encoding (SimVision)]}
    \vspace{2.2cm}}}
    \caption{Waveform mô phỏng FSM Experiment 3 -- Binary Encoding.}
    \label{fig:wave-exp3-bin}
\end{figure}

\subsubsection{Phân tích từng giai đoạn}

\paragraph{Giai đoạn 1 -- Reset đồng bộ.}
Tại chu kỳ đầu tiên với \texttt{rst = 1}, trên cạnh lên của clock, FSM được cưỡng bức về trạng thái \textbf{A} (\texttt{state = 4'b0000} với binary; \texttt{state = 9'b000000001} với one-hot). Đầu ra $z = 0$. Sau khi \texttt{rst} deasserted, FSM bắt đầu vận hành bình thường.

\paragraph{Giai đoạn 2 -- Bốn bit 0 liên tiếp ($w = 0$).}
Với $w = 0$ trong bốn chu kỳ liên tiếp, trạng thái chuyển tuần tự:
\[
A \xrightarrow{w=0} B \xrightarrow{w=0} C \xrightarrow{w=0} D \xrightarrow{w=0} E
\]
Tại chu kỳ thứ tư (trạng thái \textbf{E}), $z$ chuyển lên \textbf{1}. Đây là điểm xác nhận mạch phát hiện đúng chuỗi bốn bit 0.

\paragraph{Giai đoạn 3 -- Overlapping (tiếp tục $w = 0$).}
Khi $w = 0$ thêm một chu kỳ sau khi đạt E, FSM ở lại E nhờ self-loop (\texttt{next\_state[4] = \textasciitilde{}w \& (state[3] | state[4])}). Đầu ra $z$ tiếp tục bằng 1, xác nhận tính năng overlapping hoạt động đúng.

\paragraph{Giai đoạn 4 -- Chuyển sang đếm bit 1.}
Khi $w = 1$, FSM chuyển từ E sang \textbf{F} và $z$ trở về 0 (vì $z = \text{state}[4] | \text{state}[8]$, mà F không phải là state[4] hay state[8]). FSM bắt đầu đếm chuỗi bit 1 từ F.

\paragraph{Giai đoạn 5 -- Bốn bit 1 liên tiếp ($w = 1$).}
Trạng thái chuyển:
\[
F \xrightarrow{w=1} G \xrightarrow{w=1} H \xrightarrow{w=1} I
\]
Tại trạng thái \textbf{I}, $z = 1$. Self-loop tại I khi $w = 1$ đảm bảo overlapping cho chuỗi bit 1.

\paragraph{Giai đoạn 6 -- Chuỗi bị ngắt (mixed).}
Ví dụ: từ trạng thái B với $w = 1$, FSM chuyển sang F (không phải đếm bit 0 nữa). Kết quả là $z$ không được kích hoạt cho chuỗi bit 0 dở dang. Điều này xác nhận FSM không phát hiện sai chuỗi bị gián đoạn.

\begin{table}[H]
\centering
\caption{Tóm tắt các trường hợp điển hình trên waveform -- Experiment 3.}
\label{tab:wave-exp3}
\begin{tabular}{clccc}
\toprule
\textbf{Case} & \textbf{Mô tả} & \textbf{Trạng thái cuối} & \textbf{$z$} & \textbf{Giải thích} \\
\midrule
1 & 4 bit 0 liên tiếp         & E & 1 & Đếm đủ 4 bit 0 \\
2 & Overlapping 0 (5 bit 0)   & E & 1 & Self-loop tại E \\
3 & 4 bit 1 liên tiếp         & I & 1 & Đếm đủ 4 bit 1 \\
4 & Overlapping 1 (5 bit 1)   & I & 1 & Self-loop tại I \\
5 & Chuỗi bị ngắt (3+1)       & F & 0 & Không đạt E hay I \\
6 & Reset đồng bộ              & A & 0 & Cưỡng bức về A \\
\bottomrule
\end{tabular}
\end{table}

\begin{resultbox}
\textbf{Kết luận Experiment 3:} Cả hai hiện thực (one-hot và binary) hoạt động chính xác theo đặc tả. One-hot encoding cho thấy rõ từng flip-flop tương ứng một trạng thái trên waveform; binary encoding biểu diễn trạng thái bằng giá trị số. Cả hai cho kết quả $z$ giống nhau với cùng chuỗi $w$.
\end{resultbox}

% ============================================================
\section{Experiment 4: FSM -- Behavioral Style}
\label{sec:exp4}
% ============================================================

\subsection{Mô tả bài toán}
\label{subsec:exp4-desc}

Experiment 4 hiện thực lại \textbf{cùng FSM} của Experiment 3, nhưng sử dụng \textbf{phong cách behavioral}: logic chuyển trạng thái được mô tả trong \texttt{always\_comb} với \texttt{case} statement và \texttt{typedef enum}, thay vì viết tường minh từng phương trình Boolean.

\begin{infobox}
\textbf{Điểm khác biệt chính so với Experiment 3:}
\begin{itemize}
    \item Exp.~3 dùng \texttt{assign} với chuỗi \texttt{?:} (logic-expression) hoặc phương trình Boolean trực tiếp.
    \item Exp.~4 dùng \texttt{typedef enum} + \texttt{always\_comb} + \texttt{case} (behavioral).
    \item Cả hai mô tả cùng một phần cứng; công cụ tổng hợp tạo ra mạch tương đương.
\end{itemize}
\end{infobox}

\subsection{Cấu trúc thư mục dự án}
\label{subsec:exp4-tree}

\begin{lstlisting}[style=bash_custom, caption={Cấu trúc thư mục Experiment 4.}]
Lab1_xcelium/
|-- 00_src/
|   |-- fsm_binary_behav.sv
|   `-- fsm_onehot_behav.sv
|-- 01_tb/
|   |-- tb_fsm_binary_behav.sv
|   `-- tb_fsm_onehot_behav.sv
`-- filelist.f
\end{lstlisting}

\subsection{File \texttt{filelist.f}}
\label{subsec:exp4-filelist}

\begin{lstlisting}[style=plainfile, caption={Nội dung file \texttt{filelist.f} -- Experiment 4.}]
./00_src/fsm_binary_behav.sv
./00_src/fsm_onehot_behav.sv
./01_tb/tb_fsm_binary_behav.sv
./01_tb/tb_fsm_onehot_behav.sv
\end{lstlisting}

\subsection{File \texttt{00\_src} -- RTL Design}
\label{subsec:exp4-src}

\subsubsection{Binary Encoding: \texttt{fsm\_binary\_behav.sv}}

\begin{lstlisting}[caption={\texttt{00\_src/fsm\_binary\_behav.sv} -- Binary Encoding, Behavioral Style.},
                   label={lst:exp4-binary}]
`timescale 1ns/1ns

module fsm_binary_behav (
    input  logic clk,
    input  logic rst,
    input  logic w,
    output logic z
);

    // Binary encoding using typedef enum
    typedef enum logic [3:0] {
        A = 4'd0,
        B = 4'd1,
        C = 4'd2,
        D = 4'd3,
        E = 4'd4,
        F = 4'd5,
        G = 4'd6,
        H = 4'd7,
        I = 4'd8
    } state_t;

    state_t present_state, next_state;

    // State register (asynchronous reset)
    always_ff @(posedge clk, posedge rst) begin
        if (rst)
            present_state <= A;
        else
            present_state <= next_state;
    end

    // Next-state logic (behavioral)
    always_comb begin
        next_state = present_state;   // default: hold
        case (present_state)
            A: if (w) next_state = F; else next_state = B;
            B: if (w) next_state = F; else next_state = C;
            C: if (w) next_state = F; else next_state = D;
            D: if (w) next_state = F; else next_state = E;
            E: if (w) next_state = F; else next_state = E;
            F: if (w) next_state = G; else next_state = B;
            G: if (w) next_state = H; else next_state = B;
            H: if (w) next_state = I; else next_state = B;
            I: if (w) next_state = I; else next_state = B;
            default: next_state = A;
        endcase
    end

    // Output logic
    assign z = (present_state == E) || (present_state == I);

endmodule
\end{lstlisting}

\subsubsection{One-Hot Encoding: \texttt{fsm\_onehot\_behav.sv}}

\begin{lstlisting}[caption={\texttt{00\_src/fsm\_onehot\_behav.sv} -- One-Hot Encoding, Behavioral Style.},
                   label={lst:exp4-onehot}]
`timescale 1ns/1ns

module fsm_onehot_behav (
    input  logic clk,
    input  logic rst,
    input  logic w,
    output logic z
);

    typedef enum logic [8:0] {
        A = 9'b000000001,
        B = 9'b000000010,
        C = 9'b000000100,
        D = 9'b000001000,
        E = 9'b000010000,
        F = 9'b000100000,
        G = 9'b001000000,
        H = 9'b010000000,
        I = 9'b100000000
    } state_t;

    state_t present_state, next_state;

    // State register (asynchronous reset)
    always_ff @(posedge clk, posedge rst) begin
        if (rst)
            present_state <= A;
        else
            present_state <= next_state;
    end

    // Next-state logic (behavioral)
    always_comb begin
        next_state = present_state;   // default: hold
        case (present_state)
            A: if (w) next_state = F; else next_state = B;
            B: if (w) next_state = F; else next_state = C;
            C: if (w) next_state = F; else next_state = D;
            D: if (w) next_state = F; else next_state = E;
            E: if (w) next_state = F; else next_state = E;
            F: if (w) next_state = G; else next_state = B;
            G: if (w) next_state = H; else next_state = B;
            H: if (w) next_state = I; else next_state = B;
            I: if (w) next_state = I; else next_state = B;
            default: next_state = A;
        endcase
    end

    // Output logic
    assign z = (present_state == E) || (present_state == I);

endmodule
\end{lstlisting}

\subsection{File \texttt{01\_tb} -- Testbench}
\label{subsec:exp4-tb}

\subsubsection{Testbench Binary: \texttt{tb\_fsm\_binary\_behav.sv}}

\begin{lstlisting}[caption={\texttt{01\_tb/tb\_fsm\_binary\_behav.sv} -- Testbench cho fsm\_binary\_behav.},
                   label={lst:exp4-tb-bin}]
`timescale 1ns/1ns

module tb_fsm_binary_behav;
    logic clk;
    logic rst;
    logic w;
    logic z;

    fsm_binary_behav dut (.clk(clk), .rst(rst), .w(w), .z(z));

    initial clk = 0;
    always #5 clk = ~clk;

    // Task: apply w and check expected z value
    task apply_and_check(input logic w_in, input logic z_exp);
    begin
        @(negedge clk);
        w = w_in;
        @(posedge clk); #1;
        if (z !== z_exp)
            $display("FAIL at time=%0t: w=%0b, expected z=%0b, got z=%0b",
                     $time, w_in, z_exp, z);
        else
            $display("PASS at time=%0t: w=%0b, z=%0b",
                     $time, w_in, z);
    end
    endtask

    initial begin
        $dumpfile("tb_fsm_binary_behav.vcd");
        $dumpvars(0, tb_fsm_binary_behav);

        // Reset
        rst = 1;
        w   = 0;
        @(posedge clk); #1;
        if (z !== 0)
            $display("FAIL after reset: expected z=0, got z=%0b", z);
        else
            $display("PASS after reset");
        rst = 0;

        // --- Test Case 1: four 1s -> z=1 at 4th cycle ---
        apply_and_check(1, 0);   // A->F
        apply_and_check(1, 0);   // F->G
        apply_and_check(1, 0);   // G->H
        apply_and_check(1, 1);   // H->I  (z=1)

        // --- Overlapping 1 ---
        apply_and_check(1, 1);   // I->I  (z=1)

        // --- Transition to 0 ---
        apply_and_check(0, 0);   // I->B
        apply_and_check(0, 0);   // B->C
        apply_and_check(0, 0);   // C->D
        apply_and_check(0, 1);   // D->E  (z=1)

        // --- Overlapping 0 ---
        apply_and_check(0, 1);   // E->E  (z=1)

        $display("Method binary encoding: Da test xong");
        $finish;
    end
endmodule
\end{lstlisting}

\subsubsection{Testbench One-Hot: \texttt{tb\_fsm\_onehot\_behav.sv}}

\begin{lstlisting}[caption={\texttt{01\_tb/tb\_fsm\_onehot\_behav.sv} -- Testbench cho fsm\_onehot\_behav.},
                   label={lst:exp4-tb-oh}]
`timescale 1ns/1ns

module tb_fsm_onehot_behav;
    logic clk;
    logic rst;
    logic w;
    logic z;

    fsm_onehot_behav dut (.clk(clk), .rst(rst), .w(w), .z(z));

    // Clock generation: period = 10ns
    initial clk = 0;
    always #5 clk = ~clk;

    // Task: apply w and check expected z value
    task apply_and_check(input logic w_in, input logic z_exp);
    begin
        @(negedge clk);
        w = w_in;
        @(posedge clk); #1;
        if (z !== z_exp)
            $display("FAIL at time=%0t: w=%0b, expected z=%0b, got z=%0b",
                     $time, w_in, z_exp, z);
        else
            $display("PASS at time=%0t: w=%0b, z=%0b",
                     $time, w_in, z);
    end
    endtask

    initial begin
        $dumpfile("tb_fsm_onehot_behav.vcd");
        $dumpvars(0, tb_fsm_onehot_behav);

        // Reset
        rst = 1;
        w   = 0;
        @(posedge clk); #1;
        if (z !== 0)
            $display("FAIL: expected z=0, got z=%0b", z);
        else
            $display("PASS");
        rst = 0;

        // --- Test Case 1: four 0s -> z=1 ---
        apply_and_check(0, 0);   // A->B
        apply_and_check(0, 0);   // B->C
        apply_and_check(0, 0);   // C->D
        apply_and_check(0, 1);   // D->E  (z=1)

        // --- Overlapping 0 ---
        apply_and_check(0, 1);   // E->E  (z=1)

        // --- Test Case 2: switch to 1s ---
        apply_and_check(1, 0);   // E->F
        apply_and_check(1, 0);   // F->G
        apply_and_check(1, 0);   // G->H
        apply_and_check(1, 1);   // H->I  (z=1)

        // --- Overlapping 1 ---
        apply_and_check(1, 1);   // I->I  (z=1)

        // --- Test Case 3: broken chain ---
        apply_and_check(0, 0);   // I->B
        apply_and_check(1, 0);   // B->F
        apply_and_check(0, 0);   // F->B
        apply_and_check(0, 0);   // B->C
        apply_and_check(0, 0);   // C->D
        apply_and_check(0, 1);   // D->E  (z=1)

        $display("Method one-hot: Da test xong");
        $finish;
    end
endmodule
\end{lstlisting}

\subsection{Giải thích Waveform -- Experiment 4}
\label{subsec:exp4-wave}

% ----------------------------------------------------------
\includegraphics[width=\textwidth]{wave_exp4_binary.png}
% ----------------------------------------------------------
\begin{figure}[H]
    \centering
    \caption{Waveform mô phỏng FSM Experiment 4 -- Binary Behavioral Style.}
    \label{fig:wave-exp4-bin}
\end{figure}

% ----------------------------------------------------------
\includegraphics[width=\textwidth]{wave_exp4_onehot.png}
% ----------------------------------------------------------
\begin{figure}[H]
    \centering
    \caption{Waveform mô phỏng FSM Experiment 4 -- One-Hot Behavioral Style.}
    \label{fig:wave-exp4-oh}
\end{figure}

\subsubsection{Phân tích giai đoạn -- Binary Behavioral (\texttt{tb\_fsm\_binary\_behav})}

\paragraph{Reset không đồng bộ.}
Trong Experiment 4, reset được khai báo là \textbf{asynchronous} (\texttt{posedge rst} trong sensitivity list của \texttt{always\_ff}). Điều này có nghĩa là khi \texttt{rst = 1}, FSM chuyển về trạng thái A \textbf{ngay lập tức} mà không cần đợi cạnh clock -- khác với Experiment 3 dùng reset đồng bộ. Trên waveform, \texttt{present\_state} nhảy về A ngay khi \texttt{rst} lên 1.

\paragraph{Test Case 1 -- Bốn bit 1, z = 1.}
Testbench gọi \texttt{apply\_and\_check(1, 0)} ba lần (trạng thái F, G, H) rồi \texttt{apply\_and\_check(1, 1)} (trạng thái I, $z = 1$). Dòng \texttt{PASS} được in ra terminal, xác nhận mạch hoạt động đúng.

\paragraph{Overlapping 1.}
Sau khi đạt I, testbench gọi thêm \texttt{apply\_and\_check(1, 1)}. FSM ở lại I và $z$ vẫn là 1. Terminal in \texttt{PASS}, xác nhận self-loop tại I đúng.

\paragraph{Test Case 2 -- Chuyển sang đếm 0.}
Sau khi $w = 0$, FSM rời I về B, rồi C, D, E. Tại E, $z = 1$. Các bước trung gian (B, C, D) đều in \texttt{PASS z=0}, và bước đến E in \texttt{PASS z=1}.

\subsubsection{Phân tích giai đoạn -- One-Hot Behavioral (\texttt{tb\_fsm\_onehot\_behav})}

Testbench one-hot kiểm tra theo thứ tự ngược: đếm bốn 0 trước, rồi bốn 1, rồi mixed. Trên waveform, tín hiệu \texttt{present\_state} hiển thị dưới dạng 9-bit one-hot, ví dụ \texttt{9'b000010000} khi ở state E. Điều này giúp dễ đọc hơn trên SimVision vì có thể decode trực tiếp bit nào đang active.

\paragraph{Chuỗi bị gián đoạn (broken chain).}
Tại Test Case 3: sau \texttt{I->B}, testbench gửi $w = 1$ (B→F) rồi $w = 0$ (F→B). Việc gián đoạn giữa chừng đưa FSM về lại B, buộc phải đếm lại từ đầu. Ba bit 0 tiếp theo (B→C→D→E) và $z = 1$ tại E. Terminal in \texttt{PASS} cho tất cả các bước.

\begin{table}[H]
\centering
\caption{Tóm tắt kết quả testbench Experiment 4 -- cả hai encoding.}
\label{tab:wave-exp4}
\begin{tabular}{clccl}
\toprule
\textbf{Case} & \textbf{Mô tả} & \textbf{State cuối} & \textbf{$z$} & \textbf{Kết quả TB} \\
\midrule
Reset  & \texttt{rst=1} (async)       & A & 0 & PASS \\
1      & 4 bit 1 liên tiếp            & I & 1 & PASS \\
2      & Overlapping 1                & I & 1 & PASS \\
3      & 4 bit 0 liên tiếp            & E & 1 & PASS \\
4      & Overlapping 0                & E & 1 & PASS \\
5      & Chuỗi bị ngắt (broken chain) & E & 1 & PASS \\
\bottomrule
\end{tabular}
\end{table}

\begin{resultbox}
\textbf{Kết luận Experiment 4:} Phong cách behavioral với \texttt{typedef enum} và \texttt{case} statement cho phép mô tả FSM rõ ràng và dễ đọc hơn so với logic-expression. Cả hai phương pháp mã hóa (one-hot và binary) cho kết quả mô phỏng chính xác. Reset asynchronous đảm bảo FSM luôn trở về trạng thái xác định bất kể trạng thái clock.
\end{resultbox}

% ============================================================
\section{So sánh các phương pháp hiện thực}
\label{sec:compare}
% ============================================================

\begin{table}[H]
\centering
\caption{So sánh bốn hiện thực FSM trong Experiment 3 \& 4.}
\label{tab:compare}
\begin{tabular}{lllll}
\toprule
\textbf{Thuộc tính} & \textbf{Exp3 Binary} & \textbf{Exp3 One-Hot} & \textbf{Exp4 Binary} & \textbf{Exp4 One-Hot} \\
\midrule
Phong cách       & Logic-expr    & Logic-expr     & Behavioral      & Behavioral \\
Mã hóa state     & 4-bit binary  & 9-bit one-hot  & 4-bit binary    & 9-bit one-hot \\
Mô tả next-state & \texttt{assign ?:} & \texttt{assign [k]=} & \texttt{always\_comb case} & \texttt{always\_comb case} \\
Reset            & Synchronous   & Synchronous    & Asynchronous    & Asynchronous \\
Số flip-flop     & 4             & 9              & 4               & 9 \\
Kết quả $z$      & Giống nhau    & Giống nhau     & Giống nhau      & Giống nhau \\
\bottomrule
\end{tabular}
\end{table}

Mặc dù có sự khác biệt về cách biểu diễn và số lượng flip-flop, tất cả bốn hiện thực đều mô tả \textbf{cùng một hành vi} (behavior). Công cụ tổng hợp (synthesis tool) sẽ tối ưu hóa và ánh xạ chúng thành mạch logic tương đương tùy theo ràng buộc thiết kế.

\end{document}